\subsection{克利福德代数的结构.}

  在本节中，我们将假设 A 是一个域，E 是 A 上有限维 m 的向量空间，且二次型 Q 是非退化的（由 § 5 第 1 小节定理 1，当 A 的特征为 2 时这要求 m 为偶数）。由于 E 是自由的，典范映射 ρ 是单射（第 3 小节定理 1 的推论 2）。今后我们将把 E 与它在 C(Q) 中的象等同起来。

定理 2. --- 假设 E 的维数是偶数 \(m = 2r\)，且 Q 是中性的（\(\S 4, \text{n}^{\circ} 2\)）。则代数 C(Q) 是可分的（第八章 \(\S 7, \text{n}^{\circ} 5\)，定义 1），并且同构于 A 上维数为 \(2^{r}\) 的向量空间的自同态代数。此外，若 \(m > 0\)，则 C \(^{+}\) (Q) 是可分的，且是两个理想的直和，这两个理想都同构于 A 上维数为 \(2^{r-1}\) 的向量空间的自同态代数。

  事实上，由于 Q 是中性的，可以把 E 分解为两个维数为 r 的全奇异子空间 N 与 P 的直和（§ 4 第 2 小节命题 2 的推论 1）。由于 Q 在 N 上的限制为零，C(Q) 的由 N 生成的子代数 S 就等同于 N 的外代数（第 3 小节引理 4）。对 \(n \in \mathbb{N}\)，我们将用 \(e_n'\) 记从 S 到其自身的映射 \(t \to nt\)。

  设 \((n_{1},\ldots,n_{r})\) 是 N 的一个基；我们将用 \((p_{1},\ldots,p_{r})\) 记 P 的使 \(\Phi(n_{i},p_{j})=\delta_{ij}\) 的基（\(\S 4\)，第 2 小节，命题 2）。对 \(p\in P\)，我们将用 \(p'\) 记 N 上的线性型 \(n\to\Phi(n,p)\)，并用 \(i_{p}\) 记 S 的自同态，它是由 T(N) 的自同态 \(i_{p'}\)（与 \(p'\) 相伴，如第 2 小节引理 1 所述）通过过渡到商得到的。由 (5)，我们有

\[
e _ {n} ^ {\prime} \circ i _ {p} + i _ {p} \circ e _ {n} ^ {\prime} = \Phi (n, p) \quad (n \in \mathrm{N}, p \in \mathrm{P}).\tag{12}
\]

  对 \(x = n + p \in \mathbf{E}\)（其中 \(n \in \mathbf{N}\)，\(p \in \mathbf{P}\)），置 \(s(x) = e_n' + i_p\)。显然 \(s\) 是从 \(\mathbf{E}\) 到 \(\mathfrak{L}(\mathbf{S})\) 中的一个线性映射。由于我们有

\[
s (x) ^ {2} = \left(e _ {n} ^ {\prime} + i _ {p}\right) ^ {2} = Q (n) + \Phi (n, p) = Q (x)
\]

  根据 (12) 与引理 1（第 2 小节），s 延拓为 C(Q) 到 ℒ(S) 中的一个同态（我们仍用 s 记它）（第 1 小节命题 1）。我们将证明这个同态是满的；由于 C(Q) 与 ℒ(S) 的维数都是 2²ʳ，这将蕴涵 s 是一个同构，并证明我们的第一个断言。

  事实上，用 I 表示区间 \([1, r]\)。对 I 的每个子集 H，置 \(H' = I - H\)，并用 \(n_{H}\)（相应地，\(p_{H}\)）记诸 \(n_{i}\)（相应地，\(p_{i}\)）（\(i \in H\)，按指标递增顺序排列）的乘积。回忆 \(n_{H}\) 构成 S 的一个基（第 3 小节定理 1）。最后，对 I 的任意两个子集 H、K，置 \(x_{H,K} = n_{H}p_{I}n_{K'}\)。我们将证明元素 \(s(x_{H,K})\)（\(s(C(Q))\) 中的）生成 L(S)。现在，若 \(j \notin H\)，则由引理 1 有 \(s(p_{j})(n_{H}) = i_{p_{j}}(n_{H}) = 0\)，因为诸 \(n_{i}\)（\(i \in H\)）属于 N 上线性型 \(n \to \Phi(n, p_{j})\) 的核；另一方面我们有

\[
s \left(p _ {j}\right) \left(n _ {j} n _ {\mathrm{H}}\right) = \left(i _ {p _ {j}} \circ e _ {n _ {j}} ^ {\prime}\right) \left(n _ {\mathrm{H}}\right) = \Phi \left(p _ {j}, n _ {j}\right) n _ {\mathrm{H}} - n _ {j}. s \left(p _ {j}\right) \left(n _ {\mathrm{H}}\right) = n _ {\mathrm{H}}
\]

（由 (12)）。由于 \(s\) 是一个同态，我们对 I 的任意两个子集 H、K 推出 \(s(p_{\mathrm{K}})(n_{\mathrm{H}}) = 0\)（若 \(\mathrm{K} \not\subset \mathrm{H}\)），以及 \(s(p_{\mathbf{K}}) (n_{\mathbf{H}}) = \pm n_{\mathbf{H} - \mathbf{K}}\)（若 \(\mathrm{K} \subset \mathrm{H}\)）。由于对 \(\mathrm{M} \subset \mathrm{I}\) 与 \(\mathrm{L} \subset \mathrm{I}\)，按定义有 \(s(n_{\mathbf{M}})(n_{\mathbf{L}}) = n_{\mathbf{M}}n_{\mathbf{L}}\)，且 \(n_{\mathbf{M}}n_{\mathbf{L}}\) 在 \(\mathrm{M} \cap \mathrm{L} \neq \emptyset\) 时为零，在相反的情形等于 \(\pm n_{\mathbf{M} \cup \mathbf{L}}\)，我们由前面这些推出，对任意子集 \(\mathrm{H}\)、\(\mathrm{K}\)、\(\mathrm{L}\)（都属于 \(\mathrm{I}\)），\(s(x_{\mathbf{H},\mathbf{K}})(n_{\mathbf{L}}) = s(n_{\mathbf{H}})s(p_{\mathbf{I}})s(n_{\mathbf{K}^{\prime}})(n_{\mathbf{L}})\) 在 \(\mathrm{K} \neq \mathrm{L}\) 时为零，且等于 \(\pm n_{\mathbf{H}}\)（若 \(\mathrm{K} = \mathrm{L}\)）。这表明诸 \(s(x_{\mathbf{H},\mathbf{K}})\) 生成 \(\mathfrak{L}(\mathrm{S})\)，从而完成了第一个断言的证明。

  为证明第二个断言，置 \(S^{+} = S \cap C^{+}\) 与 \(S^{-} = S \cap C^{-}\)；显然 \(S^{+}\)（相应地，\(S^{-}\)）是 S 的由 H 的元素个数为偶数（相应地，奇数）的那些 \(n_{\mathrm{H}}\) 生成的子空间，S 是 \(S^{+}\) 与 \(S^{-}\) 的直和，且 \(s(C^{+})\) 保持 \(S^{+}\) 与 \(S^{-}\) 稳定。因此 \(s\) 把 \(C^{+}\) 映到 \(\mathfrak{L}(S)\) 的一个子代数中，它同构于乘积 \(\mathfrak{L}(S^{+}) \times \mathfrak{L}(S^{-})\)；\(s\) 在 \(C^{+}\) 上的限制是 \(C^{+}\) 到这个子代数上的一个同构，因为 \(s\) 是单射且 \(C^{+}\) 与 \(\mathfrak{L}(S^{+}) \times \mathfrak{L}(S^{-})\) 的维数都是 \(2^{2r-1}\)（第 2 小节定理 1 的推论 1）。证毕。

推论. --- 若 m 为偶数而 Q 的指标任意，则代数 C(Q) 是维数为 \(2^{m}\) 的中心单代数。此外，若 m \textgreater{} 0，子代数 C⁺(Q) 是可分的，且其中心 Z 在 A 上的维数为 2。当 Z 是一个域时，Z 是 A 的一个可分二次扩张且 C⁺(Q) 是单的；否则，Z 是同构于 A 的两个域的直积，而 C⁺(Q) 这时是维数为 \(2^{m-2}\) 的两个单子代数的直积。

  事实上，设 A' 是 A 的代数闭包，Q' 是经纯量扩张从 Q 得到的 \(E' = A' \otimes_{A} E\) 上的二次型。我们已看到 \(C(Q')\) 同构于 \(A' \otimes_{A} C(Q)\)（命题 2），且显然 \(C^{+}(Q')\) 同构于 \(A' \otimes_{A} C^{+}(Q)\)。由于 Q' 是中性的（§ 4 第 2 小节命题 3 的推论 2），本推论是定理 2 与第八章 § 7 的永存性定理的直接推论。

注. --- 1) 由于代数 C(Q) 是单的，它只有一类不可约表示；这些表示称为旋量表示；当注意力固定在这些表示之一（记作 \(\tau\)）上时，\(\tau\) 所在空间的元素称为旋量。若 Q 是中性的，则 \(\tau\) 在 \(C^{+}(Q)\) 上的限制与 \(s\) 的一样，是两个不等价绝对不可约表示之和；承载这两个表示的子空间的元素称为半旋量。在一般情形，若 \(C^{+}(Q)\) 不是单的，则 \(\tau\) 在 \(C^{+}(Q)\) 上的限制由于是忠实的，必须包含属于 \(C^{+}(Q)\) 的两类不可约表示的每一类的子表示，从而是两个不等价绝对不可约表示之和，因为在把纯量扩张到 A 的代数闭包 A' 后情况就是如此。另一方面，若 \(C^{+}(Q)\) 是单的，则它只有一类不可约表示，于是 \(\tau\) 在 \(C^{+}(Q)\) 上的限制是不可约的，因为它在把纯量扩张到 A' 后分解为两个不等价的表示。

\begin{enumerate}
\def\labelenumi{\arabic{enumi})}
\setcounter{enumi}{1}
\tightlist
\item
  假设 A 的特征 \(\neq 2\)，并设 \((x_{1},\ldots,x_{m})\)（\(m=2r\)）是 E 的一个正交基。置
\end{enumerate}

\[
z = 2 ^ {r} x _ {1} \dots x _ {m} \in \mathrm{C(Q)};
\]

  由于 \(x_{i}x_{j} + x_{j}x_{i} = 0\)（对 \(i \neq j\)），我们有 \(zx_{j} = -x_{j}z\)，这蕴涵 \(z\) 属于中心 \(Z\)（\(C^{+}(Q)\) 的）而不属于 \(A\)。我们有

\[
z ^ {2} = 2 ^ {2 r} (- 1) ^ {r} \mathrm{Q} (x _ {1}) \dots \mathrm{Q} (x _ {m}) = (- 1) ^ {r} \mathrm{D}
\]

  其中用 D 表示 \(\Phi\) 关于基 \((x_{j})\) 的判别式（参见习题 9）。

定理 3. --- 假设 E 的维数是奇数 \(m = 2r + 1\)（因而 A 的特征 \(\neq 2\)）。

\begin{enumerate}
\def\labelenumi{\alph{enumi})}
\item
  代数 C \(^{+}\) (Q) 是中心单的。若 Q 有最大指标 r，则 C \(^{+}\) (Q) 同构于 A 上维数为 2 \(^{r}\) 的向量空间的自同态代数。
\item
  代数 C(Q) 是可分的。其中心 Z 的维数为 2，且 C(Q) 同构于 Z⊗ \(_{A}\) C \(^{+}\) (Q)，因而是单的或是两个单子代数的直积。
\end{enumerate}

  事实上，设 \(x_0\) 是 E 的一个非迷向向量，F 是 \(x_0\) 的正交补；用 \(Q_1\) 表示 F 上的二次型 \(y \to -Q(x_0)Q(y)\)；显然 \(Q_1\) 是非退化的。由于 \(x_0y = -yx_0\)（对 \(y \in F\)），我们有 \((x_0y)^2 = -Q(x_0)Q(y) = Q_1(y)\)，因而映射 \(y \to x_0y\)（从 F 到 \(C^+(Q)\) 中）延拓为一个同态 \(h\)（从 \(C(Q_1)\) 到 \(C^+(Q)\) 中）（第 1 小节命题 1）。但 \(C(Q_1)\) 是单的（定理 2）且具有维数 \(2^{2r}\)（与 \(C^{+}(Q)\) 的相同）；由于 \(h(1) = 1\)，这蕴涵 \(h\) 是一个同构。此外，若 \(Q\) 的指标为 \(r\)，则可以选取 \(x_0\) 使 \(Q_1\) 也有指标 \(r\)（\(\S 4\)，第 2 小节，命题 3），这就证明了 \(a\)）。

  现在设 \((x_{1},\ldots ,x_{2r})\) 是 F 的一个正交基；置 \(z = x_0x_1\dots x_{2r}\)。立即可以验证 \(z\) 与 \(x_{j}\)（对 \(j = 0,\dots,2r\)）交换，故属于 C(Q) 的中心。设 Z 是 C(Q) 的由 1 与 \(z\) 生成的子空间；它是 C(Q) 的中心的一个子代数且是 A 的一个二次扩张，因为 \(z\) 是奇的且 \(z^2\) 等于纯量 \((-1)^r\mathrm{Q}(x_0)\ldots \mathrm{Q}(x_{2r})\)。考虑同态 \(\theta\)（从 \(\mathbf{Z}\otimes_{\mathbf{A}}\mathbf{C}^{+}(\mathbf{Q})\) 到 C(Q) 中），它由 \(\theta (u\otimes v) = uv\) 定义。由于 \(z\in \mathbb{C}^{-}\) 且可逆，映射 \(u\to zu\) 是模 \(\mathbf{C}^+\) 到 \(\mathbb{C}^{-}\) 上的一个同构，这蕴涵 \(\theta (\mathrm{Z}\otimes \mathrm{C}^{+})\) 包含 \(\mathbf{C}^{+}\) 与 \(\mathbb{C}^{-}\)，故与 C(Q) 重合。由于 \(\mathbf{Z}\otimes \mathbf{C}^{+}\) 与 C(Q) 有相同的维数 \(2^{2r + 1},\theta\) 是一个同构；这就证明了 \(b)\)，其中用到第八章 § 7 的结果。

注. --- \(\Phi\) 关于基 \((x_{j})_{(j=0,\ldots,2r)}\) 的判别式 D 等于 \(2^{2r+1}Q(x_{0})\ldots Q(x_{2r})\)。因此，Z 由 1 与奇元素 \(z' = 2^{r+1}z\) 生成，其中 \(z'^{2} = (-1)^{r}2D\)。于是代数 C(Q) 是单的当且仅当 \(2(-1)^{r}D\) 在 A 中不是平方。
% semantic-fragments: legacy-input-seam
\relax{}
